%outline:  1 what are LLMs, why they are important 2. issues to use LLMs for HPC 3. the approach we are using, 4. major findings. 

Language models (LMs) are models designed to understand and generate human language. In recent years, large language models (LLMs) %~\cite{zhao2023survey} 
trained on large amounts of text data have demonstrated stunning capabilities in various natural language processing and visualization tasks. They have also been widely used to process programming languages due to the similarities between natural languages and programming languages.  For example, GPT-4~\cite{bubeck2023sparks} shows early signs of artificial general intelligence. Based on a large language model trained on code~\cite{chen2021evaluating}, GitHub provides an AI assistant for developing software. 

Given the rise of LLMs, it is natural for researchers and developers in the high-performance computing community to start exploiting LMs for addressing various challenges in HPC, including code analysis, code generation, performance optimization, question answering, and so on. However, mainstream frameworks of LMs were originally designed to serve natural language processing. It is difficult for newcomers in HPC to quickly access HPC-specific datasets, models, and pipelines. For example, the current popular \face platform does not include dedicated pipelines for software analyses and optimizations. Another challenge is the entire field is evolving quickly, with new techniques emerging almost weekly, making it challenging for HPC users to keep up with the latest techniques and find relevant ones. Last but not least, there is a lack of standard, reproducible evaluation processes for LMs focusing on HPC-specific tasks. Therefore, it is difficult to have a fair comparison among different models for a given HPC task. 

In this paper, we propose a framework (named \ourwork) designed to serve HPC users as first-class citizens by including internal components and external APIs relevant to HPC-specific tasks. \ourwork's components include models, datasets, pipelines, and so on, while the APIs allow users to interact with these components to finish given HPC tasks. We highlight the contributions of our work as follows:
\begin{itemize}
    \item We design an extensible framework for including and exposing relevant machine learning components to facilitate the adoption of large language models for HPC-specific tasks. 
    \item The framework provides a set of APIs to facilitate essential operations, including code preprocessing, tokenization, integration with new data, and evaluation.
    %\todo{\textcolor{blue}{added in section 4}LangChain is picked to enable size limit challenge. Make sure later text mention this challenge-solution pair}    
    \item A set of pipelines have been developed to support common HPC tasks, including code similarity analysis, parallelism detection, question answering, and so on.   
    \item We provide HPC-specific datasets such as DRB-ML, \crawldata, and OMPQA to support various HPC pipelines.  
    \item Our work introduces standardized workflows and metrics to enable fair and reproducible evaluation of LLMs for HPC-specific tasks.
    \item Using three representative tasks, we demonstrated how the framework can be used to test a set of language models and generate leaderboards.
\end{itemize}

%The remainder of this paper is organized as follows. In the next section, we introduce the background of large language models and their applications in multiple domains. Section~\ref{sec-approach} presents the design of the LM4HPC framework, including its internal components and external APIs. Section~\ref{sec-evaluation} gives a preliminary results of the current implementation, using three representative HPC tasks to generate leaderboards of some popular models. Related work is presented in Section~\ref{sec-related-work}. Finally, Section~\ref{sec-conclusion} concludes the work and discusses future work. 
